\section{Objetivos}

\subsection{Objetivo general}

Desarrollar una aplicación web utilizando tecnologías emergentes que fortalezca las funcionalidades del Sistema Integral de Producción (SIP), permitiendo visualizar, proyectar y analizar información crítica relacionada con las órdenes de producción, con el fin de mejorar la precisión en la planificación, facilitar la detección temprana de riesgos y optimizar la capacidad operativa del sistema.

\subsection{Objetivos específicos}

\begin{enumerate}
    \item Analizar la arquitectura, el flujo de información y los lineamientos tecnológicos del Sistema Integral de Producción (SIP) con el propósito de asegurar una adecuada integración de los nuevos desarrollos con la plataforma existente.

    \item Diseñar e implementar componentes y vistas utilizando React.js que permitan consultar, filtrar y visualizar información relevante del proceso productivo, incluyendo proyecciones, estados y posibles riesgos asociados a las órdenes de producción.

    \item Integrar los componentes desarrollados con las fuentes de datos del sistema mediante el consumo de servicios, gestionando adecuadamente estados de carga, manejo de errores y actualización de la información.

    \item Mejorar la experiencia de usuario mediante el desarrollo de interfaces claras, estructuradas y alineadas con las necesidades operativas de planificación y control dentro del entorno productivo.

    \item Aplicar buenas prácticas de desarrollo de software, documentación y estandarización del código, con el fin de garantizar la mantenibilidad, escalabilidad y evolución futura de los módulos desarrollados dentro del SIP.
\end{enumerate}